\documentclass[12pt,a4paper,spanish,oneside]{book}
\usepackage[utf8]{inputenc}
\usepackage[spanish]{babel}
\usepackage[margin=2.5cm]{geometry}
\usepackage{xcolor}
\usepackage{fancyhdr}
\usepackage{graphicx}
\usepackage{hyperref}
\usepackage{array}
\usepackage{booktabs}
\usepackage{enumitem}
\usepackage{tikz}
\usepackage{tcolorbox}
\usepackage{amsmath}
\usepackage{amssymb}

% Configuración de encabezados y márgenes
\setlength{\headheight}{14.5pt}

% Configuración de colores
\definecolor{darkred}{HTML}{8B0000}
\definecolor{lightred}{HTML}{FFE6E6}
\definecolor{darkblue}{HTML}{003366}
\definecolor{lightblue}{HTML}{E6F2FF}

% Configuración de encabezados
\pagestyle{fancy}
\fancyhf{}
\fancyhead[R]{\textcolor{darkred}{\textbf{Manual PWA Parkinson}}}
\fancyfoot[C]{\thepage}
\renewcommand{\headrulewidth}{0.5pt}
\renewcommand{\headrule}{\color{darkred}\hrule width\headwidth height\headrulewidth \vskip-\headrulewidth}

% Comando personalizado para recuadros
\newtcolorbox{conceptobox}{colback=lightblue,colframe=darkblue,boxrule=1.5pt,
  left=10pt,right=10pt,top=10pt,bottom=10pt,fontupper=\small}

\newtcolorbox{advertenciabox}{colback=lightred,colframe=darkred,boxrule=1.5pt,
  left=10pt,right=10pt,top=10pt,bottom=10pt,fontupper=\small}

\newtcolorbox{ejemplobox}{colback={rgb:green,1;white,4},colframe=green,boxrule=1.5pt,
  left=10pt,right=10pt,top=10pt,bottom=10pt,fontupper=\small}

\title{
  \Huge \textbf{\textcolor{darkred}{Manual Completo de Uso}} \\[0.5cm]
  \huge \textbf{Plataforma Web de Asistencia Clínica} \\[0.3cm]
  \Large \textbf{Prevención de Caídas en Parkinson}
}

\author{
  Universidad Manuela Beltrán \\[0.2cm]
  Facultad de Ciencias de la Salud \\[0.2cm]
  \textit{Diseño y validación de herramienta digital} \\[0.5cm]
  \normalsize Angie Alejandra Amado Téllez
}

\date{\today}

\begin{document}

\frontmatter

% Portada
\begin{titlepage}
  \begin{center}
    \vspace*{2cm}
    
    \includegraphics[width=0.3\textwidth]{} % Aquí va el logo si existe
    
    \vspace{1cm}
    
    {\Large \textbf{UNIVERSIDAD MANUELA BELTRÁN}}
    
    {\large Facultad de Ciencias de la Salud}
    
    \vspace{2cm}
    
    \hrule
    
    \vspace{1cm}
    
    {\Huge \textbf{\textcolor{darkred}{Manual Completo de Uso}}}
    
    \vspace{0.5cm}
    
    {\huge \textbf{Prevención de Caídas en Parkinson}}
    
    \vspace{0.5cm}
    
    {\LARGE \textbf{Plataforma Web de Asistencia Clínica}}
    
    \vspace{1cm}
    
    \hrule
    
    \vspace{2cm}
    
    {\large \textbf{Descripción:} Guía exhaustiva para profesionales de fisioterapia, neurólogos y personal sanitario sobre el uso correcto de la herramienta digital de evaluación y monitoreo de riesgo de caídas en pacientes con Parkinson.}
    
    \vspace{2cm}
    
    {\normalsize Diseño y Validación: Angie Alejandra Amado Téllez}
    
    \vspace{3cm}
    
    {\normalsize \textbf{Versión:} 1.0}
    
    {\normalsize \textbf{Fecha:} \today}
    
  \end{center}
\end{titlepage}

% Página de advertencia
\newpage
\section*{Advertencia Legal}

\begin{advertenciabox}
  \textbf{RESTRICCIÓN DE USO:}
  
  Esta plataforma es un \textbf{Asistente de Decisión Clínica} para uso exclusivo de profesionales de la fisioterapia, neurólogos y personal sanitario debidamente acreditados con tarjeta profesional vigente.
  
  \textbf{NO SUSTITUYE} el criterio clínico profesional. Todos los usuarios asumen responsabilidad total sobre:
  \begin{itemize}
    \item Integridad y confidencialidad del paciente
    \item Cumplimiento de leyes de protección de datos (Ley 1581 de 2012 en Colombia)
    \item Validez clínica de los resultados interpretados
  \end{itemize}
  
  Uso exclusivamente en pacientes con diagnóstico confirmado de Parkinson en Estadios I-III (Escala Hoehn y Yahr).
\end{advertenciabox}

% Tabla de contenidos
\tableofcontents

\mainmatter

% CAPÍTULO 1: INTRODUCCIÓN
\chapter{Introducción General}

\section{Propósito de esta Guía}

Este manual ha sido diseñado para profesionales del área sanitaria que desean comprender y utilizar correctamente la Plataforma Web de Asistencia Clínica (PWA) para ``Prevención de Caídas en Parkinson''.

La guía asume \textbf{sin experiencia previa} en el sistema y proporciona:

\begin{enumerate}
  \item Explicación del contexto clínico y necesidad de la herramienta
  \item Paso a paso detallado de navegación
  \item Descripción científica de cada prueba clínica
  \item Interpretación correcta de resultados
  \item Integración en protocolos de intervención fisioterapéutica
\end{enumerate}

\section{¿Por Qué Existe Esta Herramienta?}

\subsection{Contexto Clínico}

La enfermedad de Parkinson es un trastorno neurodegenerativo progresivo que afecta la motricidad del paciente, generando:

\begin{itemize}
  \item \textbf{Bradicinesia:} Lentitud de movimientos
  \item \textbf{Rigidez:} Resistencia al movimiento pasivo
  \item \textbf{Temblor:} Oscilaciones involuntarias
  \item \textbf{Inestabilidad postural:} Pérdida de equilibrio y aumento de riesgo de caídas
\end{itemize}

Las caídas en pacientes con Parkinson representan complicaciones graves que resultan en:
\begin{itemize}
  \item Fracturas óseas
  \item Hospitalización prolongada
  \item Deterioro funcional acelerado
  \item Muerte en casos severos
\end{itemize}

\subsection{Solución Propuesta}

La PWA integra en una plataforma digital unificada:
\begin{itemize}
  \item Batería de pruebas clínicas validadas internacionalmente (TUG, BBS, FGA, UPDRS)
  \item Protocolos de monitoreo de sesión optimizados (30 minutos: 5/20/5)
  \item Estimulación rítmica sensorial (cueing) para mejora de marcha
  \item Registro automático de datos y generación de reportes
  \item Almacenamiento seguro en la nube (Supabase)
\end{itemize}

\section{Beneficios Clave}

\begin{conceptobox}
  \textbf{Para el Profesional:}
  \begin{itemize}
    \item Evaluación estandarizada y reproducible
    \item Reducción de tiempo de documentación manual
    \item Acceso a interpretaciones clínicas automáticas
    \item Seguimiento longitudinal del paciente
  \end{itemize}
  
  \textbf{Para el Paciente:}
  \begin{itemize}
    \item Intervención personalizada basada en datos objetivos
    \item Cueing auditivo para sincronización de marcha
    \item Sesiones de duración controlada y científicamente diseñadas
    \item Feedback en tiempo real de riesgo
  \end{itemize}
\end{conceptobox}

% CAPÍTULO 2: INFORMACIÓN GENERAL DE LA PWA
\chapter{Información General de la Plataforma}

\section{Acceso y Autenticación}

\subsection{URL de Acceso}

La plataforma está disponible en producción en:

\begin{center}
  \fcolorbox{darkred}{lightred}{
    \parbox{0.8\textwidth}{
      \centering
      \texttt{\Large https://cartilla-parkinson-angie.vercel.app}
    }
  }
\end{center}

\subsection{Escudo de Acceso Profesional}

Al acceder a la PWA, la primera pantalla presenta un \textbf{Escudo de Acceso Profesional}:

\begin{advertenciabox}
  \textbf{SISTEMA DE USO RESTRINGIDO}
  
  Antes de continuar, debe declarar:
  \begin{enumerate}
    \item Posesión de tarjeta profesional vigente (Colegio de Fisioterapeutas, Colegio Médico, etc.)
    \item Aceptación de responsabilidad sobre integridad del paciente
    \item Compromiso con confidencialidad de datos bajo Ley 1581 de 2012
  \end{enumerate}
  
  \textbf{Al hacer clic en ``Declaro ser profesional y acepto los términos'', usted asume estas responsabilidades legales.}
\end{advertenciabox}

\section{Arquitectura General de la PWA}

La plataforma se estructura en 5 módulos principales:

\begin{enumerate}
  \item \textbf{Perfil del Paciente:} Registro de datos demográficos y clínicos
  \item \textbf{Batería de Evaluación Clínica:} 4 pruebas estandarizadas
  \item \textbf{Estimulación Rítmica Sensorial (Cueing):} Metrónomo clínico
  \item \textbf{Monitoreo de Sesión:} Contador de tiempo y fases protocolizadas
  \item \textbf{Seguimiento Meta Terapéutica:} Registro de hallazgos y guardado de datos
\end{enumerate}

% CAPÍTULO 3: GUÍA DE NAVEGACIÓN PASO A PASO
\chapter{Guía Detallada de Navegación}

\section{Paso 1: Aceptación del Escudo Profesional}

\textbf{Acción:} Haga clic en el botón \texttt{``Declaro ser profesional y acepto los términos''}.

\textbf{Resultado:} La pantalla se actualiza y muestra el contenido de la PWA.

\begin{conceptobox}
  \textbf{Nota importante:} Este paso verifica que solo profesionales acreditados accedan. Si no es profesional sanitario o no tiene responsabilidad clínica sobre pacientes, \textbf{no debe continuar}.
\end{conceptobox}

\section{Paso 2: Completar Perfil del Paciente}

\subsection{Campos Obligatorios}

Una vez dentro de la PWA, la primera sección es \textbf{``Perfil del Paciente''} con los siguientes campos:

\begin{table}[h]
  \centering
  \begin{tabular}{|p{3cm}|p{4cm}|p{4cm}|}
    \hline
    \textbf{Campo} & \textbf{Tipo} & \textbf{Descripción} \\
    \hline
    Nombre Completo & Texto & Identificación del paciente \\
    \hline
    Estadio Hoehn y Yahr & Selector (I-III) & Fase de Parkinson \\
    \hline
    MMSE & Número (0-30) & Mini Mental State Examination \\
    \hline
    Profesional Responsable & Texto & Nombre del fisioterapeuta/médico \\
    \hline
    Evaluación Socio-Ambiental & Texto largo & Descripción de barreras en el hogar \\
    \hline
  \end{tabular}
\end{table}

\subsection{Explicación Detallada de Cada Campo}

\subsubsection{Nombre Completo}

\textbf{Uso:} Identificación unívoca del paciente en registros.

\textbf{Consideraciones:} 
\begin{itemize}
  \item Completar con nombre y apellidos como aparecen en documento de identidad
  \item Es crítico para garantizar trazabilidad clínica
\end{itemize}

\subsubsection{Estadio Hoehn y Yahr}

\textbf{Escala de Clasificación del Parkinson:}

\begin{conceptobox}
  \begin{itemize}
    \item \textbf{Estadio I:} Síntomas unilaterales leves, sin afectación funcional significativa
    \item \textbf{Estadio II:} Síntomas bilaterales o de línea media, sin trastornos del equilibrio
    \item \textbf{Estadio III:} Enfermedad bilateral leve, inestabilidad postural incipiente pero independencia en actividades diarias
  \end{itemize}
  
  \textbf{Nota:} Esta PWA está diseñada SOLO para Estadios I-III. Estadios IV-V requieren intervención diferente.
\end{conceptobox}

\subsubsection{MMSE (Mini Mental State Examination)}

\textbf{Descripción:} Evaluación cognitiva breve (puntuación 0-30).

\textbf{Interpretación:}
\begin{itemize}
  \item $\geq 24$: Cognición preservada (requiere esta PWA)
  \item 18-23: Deterioro cognitivo leve
  \item $<18$: Deterioro cognitivo moderado a severo
\end{itemize}

\textbf{Requisito PWA:} La plataforma requiere MMSE $\geq 24$ para garantizar capacidad de consentimiento y comprensión del protocolo.

\subsubsection{Profesional Responsable}

\textbf{Uso:} Documentación de responsabilidad legal y clínica.

\textbf{Formato sugerido:} ``Lic. [Nombre] [Apellido]'' o ``Dr./Dra. [Nombre] [Apellido]''

\subsubsection{Evaluación Socio-Ambiental}

\textbf{Descripción:} Campo de texto libre para documentar:

\begin{itemize}
  \item \textbf{Red de apoyo:} ¿Con quién vive? ¿Familiares disponibles?
  \item \textbf{Barreras en el hogar:} Alfombras sueltas, escaleras sin barandilla, iluminación deficiente
  \item \textbf{Factores de riesgo ambiental:} Pisos resbaladizos, muebles bajos
  \item \textbf{Accesibilidad:} ¿Hay espacios seguros para caminar?
\end{itemize}

\textbf{Ejemplo:}
\begin{ejemplobox}
  ``Paciente vive con esposa. Hogar de dos plantas con escaleras sin barandilla. Piso principal con alfombra en sala (factor de riesgo). Baño con barra de apoyo. Pasillos amplios y bien iluminados. Riesgo moderado por escaleras y alfombra.''
\end{ejemplobox}

\section{Paso 3: Seleccionar y Ejecutar Pruebas Clínicas}

Después de completar el perfil, el siguiente módulo es \textbf{``Batería de Evaluación Clínica''} con 4 pruebas:

\begin{center}
  \begin{tabular}{|c|c|c|}
    \hline
    \textbf{Prueba} & \textbf{Sigla} & \textbf{Función} \\
    \hline
    Timed Up and Go Cognitivo & TUG & Movilidad y doble tarea \\
    \hline
    Berg Balance Scale & BBS & Equilibrio estático/dinámico \\
    \hline
    Functional Gait Assessment & FGA & Estabilidad de marcha \\
    \hline
    UPDRS Motor & UPDRS & Progresión de síntomas motores \\
    \hline
  \end{tabular}
\end{center}

\textbf{Requisito mínimo:} Completar al menos 2 pruebas para habilitar el botón ``Guardar Valoración''.

% CAPÍTULO 4: DESCRIPCIÓN DETALLADA DE PRUEBAS
\chapter{Descripción Exhaustiva de Cada Prueba Clínica}

\section{Prueba 1: TUG Cognitivo (Timed Up and Go)}

\subsection{¿Qué es?}

El \textbf{TUG Cognitivo} es una adaptación de la prueba clásica ``Timed Up and Go'' que añade un componente de doble tarea cognitiva para evaluar simultáneamente movilidad y función ejecutiva.

\subsection{Fundamento Clínico}

En Parkinson, la doble tarea (realizar dos actividades cognitivas simultáneamente, como contar regresivamente mientras se camina) evidencia:

\begin{itemize}
  \item \textbf{Automaticidad perdida:} Pacientes normales mantienen ritmo; pacientes con Parkinson disminuyen velocidad
  \item \textbf{Fallos de atención dividida:} Manifestación de disfunción de ganglios basales
  \item \textbf{Riesgo aumentado de caídas:} Alteración de automatismo = pérdida de equilibrio
\end{itemize}

\subsection{Protocolo Detallado}

\subsubsection{Preparación}

\begin{enumerate}
  \item Paciente sentado en silla con brazos (altura estándar $\approx 46$cm)
  \item Distancia de 3 metros marcada claramente en el piso
  \item Espacio libre de obstáculos
  \item Paciente con calzado seguro (no calcetines solos)
\end{enumerate}

\subsubsection{Ejecución}

\begin{enumerate}
  \item \textbf{Instrucción al paciente:} ``Cuando diga 'ya', cuente regresivamente desde 100 de 7 en 7 (100, 93, 86...). Luego, levántese, camine 3 metros, gire 180 grados y regrese. Continúe contando todo el tiempo''.
  \item \textbf{Acción del profesional:}
  \begin{itemize}
    \item Inicia cronómetro en el sitio.io (el botón ``INICIAR'' en la PWA)
    \item Observa: bradicinesia, congelamiento, desequilibrio
    \item Nota comportamiento: ¿dejó de contar? ¿redujo velocidad de marcha?
  \end{itemize}
  \item \textbf{Fin:} Cuando paciente regresa a la silla y se sienta, presiona ``DETENER''
\end{enumerate}

\subsection{Interpretación de Resultados}

\begin{table}[h]
  \centering
  \begin{tabular}{|c|c|c|}
    \hline
    \textbf{Tiempo} & \textbf{Clasificación} & \textbf{Riesgo} \\
    \hline
    $< 10$ segundos & Normal & Bajo \\
    \hline
    11-20 segundos & Riesgo leve & Medio \\
    \hline
    $> 20$ segundos & Riesgo alto & Alto \\
    \hline
  \end{tabular}
\end{table}

\textbf{Interpretación en la PWA:}

El sistema automáticamente mostrará:
\begin{itemize}
  \item \textbf{Estado:} ``Completa''
  \item \textbf{Tiempo capturado:} (ej: ``8.5 segundos'')
  \item \textbf{Clasificación de riesgo:} ✅ Riesgo Bajo
\end{itemize}

\subsection{Consideraciones Clínicas Avanzadas}

\begin{advertenciabox}
  \textbf{Interpretación Clínica Crítica:}
  
  Si el paciente \textbf{dejó de contar regresivamente}, eso es más significativo que el tiempo total. Indica fallo atencional grave, sugestivo de:
  \begin{itemize}
    \item Déficit ejecutivo severo
    \item Posible deterioro cognitivo no diagnosticado
    \item Necesidad de valoración neuropsicológica
  \end{itemize}
  
  La PWA no captura esto automáticamente; \textbf{el profesional debe registrarlo} en el campo de observaciones.
\end{advertenciabox}

\section{Prueba 2: BBS (Berg Balance Scale)}

\subsection{¿Qué es?}

La \textbf{Berg Balance Scale} es un instrumento de 14 ítems que evalúa equilibrio estático y dinámico en pacientes con riesgo de caída. Es el estándar de oro en neurología para predicción de caídas.

\subsection{Componentes Evaluados}

\begin{table}[h]
  \centering
  \begin{tabular}{|p{5cm}|p{5cm}|}
    \hline
    \textbf{Equilibrio Estático} & \textbf{Equilibrio Dinámico} \\
    \hline
    1. Sedente a bipedestación & 4. Bipedestación a sedente \\
    2. Bipedestación sin apoyo & 5. Transferencias \\
    3. Sedente sin apoyo & 10. Giro para mirar atrás \\
    6. Bipedestación con ojos cerrados & 11. Giro 360 grados \\
    7. Bipedestación con pies juntos & 13. Bipedestación en tándem \\
    & 14. Bipedestación a una pierna \\
    \hline
  \end{tabular}
\end{table}

\subsection{Protocolo Detallado para Cada Ítem}

\subsubsection{Ítem 1: Sedente a Bipedestación}

\textbf{Instrucción:} ``Levántese de la silla sin usar brazos''

\textbf{Puntuación (0-4):}
\begin{itemize}
  \item \textbf{0:} Incapaz sin ayuda
  \item \textbf{1:} Necesita máxima asistencia
  \item \textbf{2:} Necesita asistencia mínima o supervisión
  \item \textbf{3:} Capaz sin ayuda, pero usa brazos
  \item \textbf{4:} Capaz sin ayuda, sin usar brazos
\end{itemize}

\textbf{Observación clínica en Parkinson:}
\begin{itemize}
  \item Lentitud de movimiento (bradicinesia)
  \item Pérdida de movimientos asociados (brazos)
  \item Dificultad inicial (congelamiento)
\end{itemize}

\subsubsection{Ítem 2: Bipedestación sin Apoyo}

\textbf{Instrucción:} ``Manténgase de pie sin aferrarse, durante 2 minutos''

\textbf{Puntuación (0-4):}
\begin{itemize}
  \item \textbf{0:} Requiere sujeción constante
  \item \textbf{1:} Sujeción activa continua
  \item \textbf{2:} Sujeción activa intermitente
  \item \textbf{3:} Puede mantenerse 2 minutos con dificultad
  \item \textbf{4:} Bipedestación estable sin dificultad
\end{itemize}

\subsubsection{Ítem 3: Sedente sin Apoyo}

\textbf{Instrucción:} ``Siéntese con manos cruzadas sobre pecho, durante 2 minutos''

\textbf{Evaluación:} Estabilidad del tronco en ausencia de apoyo de extremidades

\subsubsection{Ítem 4: Bipedestación a Sedente}

\textbf{Instrucción:} ``Siéntese de forma controlada''

\textbf{Observación en Parkinson:}
\begin{itemize}
  \item ¿Cae hacia atrás sin control?
  \item ¿Mira atrás antes de sentarse? (reflejo protector normal)
  \item ¿Pérdida de movimientos asociados?
\end{itemize}

\subsubsection{Ítem 5: Transferencias}

\textbf{Instrucción:} ``Muévase de una silla a otra sin armadura''

\textbf{Criterio:} Cambio de posición sin necesidad de brazos

\subsubsection{Ítem 6: Bipedestación con Ojos Cerrados}

\textbf{Instrucción:} ``Cierre los ojos y manténgase de pie durante 10 segundos''

\textbf{Crítico en Parkinson:} La propriocepción está afectada; sin visión, inestabilidad aumenta significativamente

\subsubsection{Ítem 7: Bipedestación con Pies Juntos}

\textbf{Instrucción:} ``Pies juntos, talones casi tocándose, manténgase sin apoyo''

\textbf{Base de sustentación:} Se reduce; desafío mayor de equilibrio

\subsubsection{Ítem 8: Alcance hacia Adelante con Brazo Extendido}

\textbf{Instrucción:} ``Extienda el brazo y alcance lo máximo que pueda hacia adelante sin perder equilibrio''

\textbf{Medida:} Distancia en centímetros (distancia normal: $\geq 25$cm)

\subsubsection{Ítem 9: Recoger Objeto del Piso}

\textbf{Instrucción:} ``Recoja este objeto del piso''

\textbf{Observación:} 
\begin{itemize}
  \item Técnica de flexión (cadera vs. espalda)
  \item Pérdida de equilibrio durante el descenso
\end{itemize}

\subsubsection{Ítem 10: Giro para Mirar Atrás}

\textbf{Instrucción:} ``Gire su cabeza y mire hacia atrás por encima de cada hombro''

\textbf{Riesgo en Parkinson:} Pérdida de estabilidad, especialmente en pacientes con freezing

\subsubsection{Ítem 11: Giro de 360 Grados}

\textbf{Instrucción:} ``Gire lentamente 360 grados''

\textbf{Observación:}
\begin{itemize}
  \item Velocidad (bradicinesia)
  \item Número de pasos (normal: $\leq 12$)
  \item Congelamiento durante la rotación (freezing)
\end{itemize}

\subsubsection{Ítem 12: Colocar Pie sobre Escalón}

\textbf{Instrucción:} ``Coloque un pie sobre el escalón, manteniendo equilibrio sin apoyo''

\textbf{Evaluación:} Capacidad de mantener equilibrio unipodal

\subsubsection{Ítem 13: Bipedestación en Tándem}

\textbf{Instrucción:} ``Ponga un pie directamente frente al otro (en línea), manténgase sin apoyo durante 30 segundos''

\textbf{Dificultad progresiva:} Mayor desafío que pies separados

\subsubsection{Ítem 14: Bipedestación a Una Pierna}

\textbf{Instrucción:} ``Levante una pierna sin tocar el piso, manténgase durante 10 segundos''

\textbf{Criterio de éxito:} Capacidad de mantener la posición sin perder equilibrio

\subsection{Puntuación Total BBS}

\textbf{Rango:} 0-56 puntos

\textbf{Interpretación:}

\begin{table}[h]
  \centering
  \begin{tabular}{|c|c|c|}
    \hline
    \textbf{Puntuación} & \textbf{Riesgo de Caída} & \textbf{Acción Recomendada} \\
    \hline
    46-56 & Bajo riesgo & Continuar actividad normal \\
    \hline
    40-45 & Riesgo moderado & Intervención preventiva moderada \\
    \hline
    $< 45$ & \textbf{RIESGO ALTO} & \textbf{Intervención inmediata requerida} \\
    \hline
  \end{tabular}
\end{table}

\textbf{Advertencia Clínica:}

\begin{advertenciabox}
  Una puntuación \textbf{$< 45/56$ es INDICADOR CRÍTICO de riesgo inminente de caídas múltiples}. Requiere:
  \begin{itemize}
    \item Intervención preventiva inmediata
    \item Aumento de frecuencia de seguimiento
    \item Posible ajuste de medicación
    \item Modificación ambiental del hogar
    \item Consideración de asistencia en la marcha (bastón, andador)
  \end{itemize}
\end{advertenciabox}

\subsection{Cómo Usar BBS en la PWA}

La PWA presenta 14 botones con opciones 0-4 para cada ítem:

\begin{enumerate}
  \item Ejecute cada maniobra según el protocolo descrito
  \item Asigne puntuación inmediatamente (0-4)
  \item Haga clic en el botón correspondiente
  \item La PWA suma automáticamente
  \item Al completar los 14 ítems, verá:
  \begin{itemize}
    \item Puntuación total (X/56)
    \item Interpretación automática de riesgo
    \item El estado cambia a ``Completa''
  \end{itemize}
\end{enumerate}

\section{Prueba 3: FGA (Functional Gait Assessment)}

\subsection{¿Qué es?}

\textbf{FGA} es una escala de 10 ítems que evalúa la estabilidad de la marcha bajo diferentes condiciones de desafío. Es particularmente útil en Parkinson para detectar freezing y adaptación de marcha a obstáculos.

\subsection{Los 10 Ítems}

\begin{enumerate}
  \item \textbf{Marcha en superficie plana}
  \item \textbf{Cambio de velocidad de marcha}
  \item \textbf{Marcha con giro horizontal de cabeza}
  \item \textbf{Marcha con giro vertical de cabeza}
  \item \textbf{Giro y pivote durante la marcha}
  \item \textbf{Paso sobre obstáculo}
  \item \textbf{Marcha con base estrecha}
  \item \textbf{Marcha con ojos cerrados}
  \item \textbf{Marcha hacia atrás}
  \item \textbf{Subir y bajar escaleras}
\end{enumerate}

\subsection{Protocolo y Puntuación}

Cada ítem se puntúa de 0-3:

\begin{itemize}
  \item \textbf{0:} Pérdida normal del equilibrio, puede caer (requerería sujeción inmediata)
  \item \textbf{1:} Pérdida de equilibrio, se recupera sin sujeción; puede tambalearse pero se estabiliza
  \item \textbf{2:} Inestabilidad mínima; camina con confianza
  \item \textbf{3:} Marcha normal, sin inestabilidad observable
\end{itemize}

\subsubsection{Ítem 1: Marcha en Superficie Plana}

\textbf{Instrucción:} ``Camine 10 metros a paso normal''

\textbf{Observación:}
\begin{itemize}
  \item Longitud de paso (reducida en Parkinson)
  \item Cadencia (velocidad de pasos)
  \item Balanceo de brazos (reducido)
  \item Congelamiento inicial o terminal
\end{itemize}

\subsubsection{Ítem 2: Cambio de Velocidad}

\textbf{Instrucción:} ``Camine a paso lento 5 metros, luego a paso normal 5 metros, luego rápido 5 metros''

\textbf{Crítico en Parkinson:} Dificultad para modular velocidad es característica. Pueden no lograrlo.

\subsubsection{Ítem 3 y 4: Marcha con Giros de Cabeza}

\textbf{Instrucción:} ``Camine 10 metros mientras gira la cabeza lado a lado (horizontal) / arriba y abajo (vertical)''

\textbf{Fundamento:} Los giros de cabeza durante la marcha requieren integración vestibular-cerebelosa. En Parkinson, esto causa inestabilidad.

\subsubsection{Ítem 5: Giro y Pivote durante la Marcha}

\textbf{Instrucción:} ``Camine 10 metros, luego gire 180 grados''

\textbf{Observación:}
\begin{itemize}
  \item Velocidad de rotación
  \item Congelamiento durante el giro (freezing)
  \item Base de sustentación durante la rotación
\end{itemize}

\subsubsection{Ítem 6: Paso sobre Obstáculo}

\textbf{Instrucción:} ``Camine sobre estos obstáculos sin modificar ritmo''

\textbf{Desafío:} Requiere adaptación visomotora. En Parkinson, pueden frotar con los pies.

\subsubsection{Ítem 7: Marcha con Base Estrecha}

\textbf{Instrucción:} ``Camine en línea recta (un pie frente al otro) durante 10 metros''

\textbf{Dificultad aumentada:} Reduce base de sustentación. Inestabilidad esperada en Parkinson.

\subsubsection{Ítem 8: Marcha con Ojos Cerrados}

\textbf{Instrucción:} ``Cierre los ojos y camine 10 metros (supervisor muy cercano por seguridad)''

\textbf{Crítico:} Sin aferencia visual, la propriocepción deteriorada en Parkinson causa riesgo alto de caída.

\subsubsection{Ítem 9: Marcha Hacia Atrás}

\textbf{Instrucción:} ``Camine hacia atrás 10 metros (supervisor vigilando)''

\textbf{Dificultad:} Requiere integración compleja. Riesgo muy alto en Parkinson.

\subsubsection{Ítem 10: Escaleras}

\textbf{Instrucción:} ``Suba 6-8 escalones con barandilla. Baje''

\textbf{Observación:}
\begin{itemize}
  \item Patrón de marcha (paso a paso vs. alternado)
  \item Uso de barandilla
  \item Congelamiento en inicio o al cambiar de dirección
\end{itemize}

\subsection{Puntuación Total FGA}

\textbf{Rango:} 0-30 puntos

\textbf{Interpretación:}

\begin{table}[h]
  \centering
  \begin{tabular}{|c|c|}
    \hline
    \textbf{Puntuación} & \textbf{Interpretación} \\
    \hline
    $\geq 22$ & Bajo riesgo de caída \\
    \hline
    19-21 & Riesgo moderado \\
    \hline
    $< 19$ & Alto riesgo de caída \\
    \hline
  \end{tabular}
\end{table}

\section{Prueba 4: UPDRS Motor (Unified Parkinson's Disease Rating Scale)}

\subsection{¿Qué es?}

\textbf{UPDRS} es la escala más utilizada internacionalmente para evaluar severidad de síntomas motores en Parkinson. La sección Motor (Parte III) es de particular interés en esta PWA.

\subsection{Componentes del UPDRS Motor}

\begin{table}[h]
  \centering
  \begin{tabular}{|p{6cm}|p{4cm}|}
    \hline
    \textbf{Ítem} & \textbf{Puntuación} \\
    \hline
    Rigidez & 0-4 \\
    \hline
    Bradicinesia & 0-4 \\
    \hline
    Tremor & 0-4 \\
    \hline
    Postura & 0-2 \\
    \hline
    Marcha & 0-2 \\
    \hline
    Estabilidad Postural & 0-2 \\
    \hline
  \end{tabular}
\end{table}

\subsection{En la PWA: Campo de Texto Libre}

Diferente a BBS y FGA, UPDRS en esta PWA se captura mediante un \textbf{campo de texto libre} donde el profesional registra:

\begin{itemize}
  \item Puntuación motor UPDRS parcial o total
  \item Hallazgos principales
  \item Severidad de síntomas
  \item Cambios desde última evaluación
\end{itemize}

\textbf{Ejemplo de Registro:}

\begin{ejemplobox}
  ``UPDRS motor: 18/108. Bradicinesia leve en extremidades superiores. Rigidez de rueda dentada en cuello. Tremor en reposo mínimo. Postura levemente inclinada. Marcha con pasos cortos sin congelamiento en esta visita. Estabilidad postural preservada.''
\end{ejemplobox}

\subsection{Interpretación Clínica}

\textbf{Severidad según UPDRS Motor Total:}

\begin{table}[h]
  \centering
  \begin{tabular}{|c|c|}
    \hline
    \textbf{Puntuación} & \textbf{Interpretación} \\
    \hline
    0-20 & Síntomas leves \\
    \hline
    21-50 & Síntomas moderados \\
    \hline
    51-108 & Síntomas severos \\
    \hline
  \end{tabular}
\end{table}

% CAPÍTULO 5: MÓDULOS COMPLEMENTARIOS
\chapter{Módulos Complementarios de la PWA}

\section{Estimulación Rítmica Sensorial (Cueing)}

\subsection{¿Qué es el Cueing?}

\textbf{Cueing} (del inglés: ``pista rítmica'') es una técnica neurorrehabilitadora que usa estímulos auditivos, visuales o táctiles para sincronizar la marcha del paciente.

\subsection{Fundamento Neurocientífico}

En Parkinson, los ganglios basales están dañados. Estos normalmente generan ritmo automático de movimiento. Sin ellos:

\begin{itemize}
  \item La marcha se vuelve errática
  \item Aparece \textbf{congelamiento} (freezing of gait)
  \item El paciente debe usar conscientemente la corteza cerebral para planificar cada paso
\end{itemize}

El \textbf{cueing auditivo} ``secuestra'' la corteza auditiva, permitiendo que otro sistema neural genere ritmo automático. Paradójicamente, el paciente camina mejor \textbf{siguiendo un metrónomo que caminando libremente}.

\subsection{Implementación en la PWA}

La PWA incluye un \textbf{metrónomo clínico} ajustable:

\begin{itemize}
  \item \textbf{Rango de BPM:} 40-120 (latidos por minuto)
  \item \textbf{Estímulo:} Sonido de tono puro (880 Hz, duración breve)
  \item \textbf{Indicador visual:} Círculo blanco que pulsea sincronizadamente
\end{itemize}

\subsection{Cómo Usarlo}

\begin{enumerate}
  \item \textbf{Establecer BPM:} Mueva el slider a la velocidad deseada (ej: 60 BPM)
  
  \textbf{Consideración:} Cadencia normal de marcha en adultos = 100-120 pasos/min. En Parkinson, empezar más lento (60-80 BPM)
  
  \item \textbf{Indicador visual:} Observará un círculo que se ilumina con cada pulso
  
  \item \textbf{Instrucción al paciente:} ``Camine sincronizando sus pasos con cada sonido que oye''
  
  \item \textbf{Observación:} Monitorear:
  \begin{itemize}
    \item ¿Sincronización con metrónomo?
    \item ¿Mejora de velocidad respecto a marcha libre?
    \item ¿Reducción de congelamiento?
    \item ¿Mejora de amplitud de paso?
  \end{itemize}
  
  \item \textbf{Registrar cambios:} Si hubo mejora significativa, documéntelo en ``Seguimiento Meta Terapéutica''
\end{enumerate}

\section{Monitoreo de Sesión}

\subsection{Estructura de la Sesión}

\textbf{PROTOCOLO OPTIMIZADO: 30 MINUTOS TOTALES}

La PWA implementa un protocolo de sesión de 30 minutos dividido en 3 fases:

\begin{table}[h]
  \centering
  \begin{tabular}{|c|c|c|c|}
    \hline
    \textbf{Fase} & \textbf{Duración} & \textbf{Objetivo} & \textbf{Actividades} \\
    \hline
    1. Calentamiento & 5 min & Preparación neuromuscular & Movilidad articular, respiración \\
    \hline
    2. Intervención Digital & 20 min & Aplicación de estímulos & Cueing, evaluación dinámica \\
    \hline
    3. Enfriamiento & 5 min & Retorno a homeostasis & Estiramientos, retroalimentación \\
    \hline
  \end{tabular}
\end{table}

\subsection{¿Por Qué 30 Minutos?}

Esta duración fue seleccionada basándose en evidencia clínica:

\begin{conceptobox}
  \begin{itemize}
    \item \textbf{$< 20$ min:} Insuficiente para neuroplasticidad
    \item \textbf{20-30 min:} Óptimo para sesión de intervención (tolerancia neuromuscular, atención)
    \item \textbf{$> 45$ min:} Fatiga cognitiva, fatiga motora, reducción de atención
  \end{itemize}
  
  El protocolo 5/20/5 respeta los principios de rehabilitación neuromuscular mientras minimiza fatiga.
\end{conceptobox}

\subsection{Uso del Cronómetro en la PWA}

\begin{enumerate}
  \item \textbf{Inicio:} Haga clic en el botón \texttt{``Play''} (triángulo) al comenzar intervención
  
  \item \textbf{Visualización:} La PWA muestra tiempo regresivo (30:00 → 00:00)
  
  \item \textbf{Cambio de fases:} Los bloques de ``Calentamiento'', ``Intervención Digital'', ``Enfriamiento'' se iluminan automáticamente según tiempo restante:
  \begin{itemize}
    \item Bloque 1 iluminado: segundos $> 1500$ (25+ minutos)
    \item Bloque 2 iluminado: segundos entre 1500 y 300 (5-25 min)
    \item Bloque 3 iluminado: segundos $< 300$ (últimos 5 minutos)
  \end{itemize}
  
  \item \textbf{Pausa:} Haga clic en ``Pausa'' (dos líneas) si necesita interrumpir
  
  \item \textbf{Reinicio:} Botón ``Reinicio'' resetea a 30:00
\end{enumerate}

\subsection{Monitoreo Clínico Simultáneo}

\textbf{El profesional debe:}

\begin{enumerate}
  \item Observar respuesta del paciente a cada estímulo
  \item Anotar cambios en:
  \begin{itemize}
    \item Fatiga (escala 1-10)
    \item Calidad de marcha
    \item Participación/motivación
    \item Eventos adversos (caídas, mareos, dolor)
  \end{itemize}
  \item Documentar en ``Seguimiento Meta Terapéutica''
\end{enumerate}

\section{Seguimiento Meta Terapéutica}

\subsection{Campo de Registro Diario}

Este módulo final permite documentación de:

\begin{itemize}
  \item \textbf{Observaciones clínicas:} Hallazgos principales de la sesión
  \item \textbf{Respuesta del paciente:} Tolerancia, fatiga, motivación
  \item \textbf{Cambios funcionales:} Mejora o deterioro en pruebas
  \item \textbf{Recomendaciones:} Ajustes para próxima sesión
\end{itemize}

\subsection{Requisito de Integridad}

\textbf{Mínimo requerido para guardar:}

\begin{advertenciabox}
  \textbf{Completar al menos 2 pruebas clínicas (de 4 disponibles)}
  
  Esto garantiza:
  \begin{itemize}
    \item Evaluación multimodal del paciente
    \item Confiabilidad de interpretación clínica
    \item Cumplimiento de estándares de protocolización
  \end{itemize}
  
  Hasta que no cumpla este requisito, el botón ``Guardar Valoración'' permanecerá deshabilitado.
\end{advertenciabox}

\subsection{Ejemplo de Registro Completo}

\begin{ejemplobox}
  \textbf{Paciente:} Juan Pérez López \\
  \textbf{Estadio:} II \\
  \textbf{Fecha:} 20/05/2026 \\
  \textbf{Profesional:} Lic. Laura Gómez
  
  \vspace{0.3cm}
  
  \textbf{Pruebas Realizadas:}
  \begin{itemize}
    \item TUG: 8.5 s (Completa - Riesgo Bajo)
    \item UPDRS Motor: 18/108 (Síntomas leves)
  \end{itemize}
  
  \textbf{Observaciones:}
  ``Sesión tolerada adecuadamente. Paciente refiere mejoría subjetiva con cueing auditivo a 70 BPM. Marcha más fluida, menos ``pasos arrastrados''. Bradicinesia leve en extremidad superior derecha pero sin cambios desde sesión anterior. Paciente motivado y adherente al protocolo. Tolerancia completa de 30 minutos sin fatiga.''
  
  \textbf{Recomendaciones:}
  ``Continuar con cueing a 70 BPM. Incrementar próxima sesión a 75 BPM si tolera. Realizar seguimiento en 1 semana. Reforzar ejercicios en casa (marcha con metrónomo 10 min diarios).''
\end{ejemplobox}

% CAPÍTULO 6: CASOS CLÍNICOS PRÁCTICOS
\chapter{Casos Clínicos Prácticos}

\section{Caso 1: Paciente con Riesgo Bajo}

\subsection{Presentación}

\begin{itemize}
  \item \textbf{Nombre:} Carlos Méndez Rojas
  \item \textbf{Edad:} 68 años
  \item \textbf{Estadio Hoehn-Yahr:} II
  \item \textbf{Diagnóstico:} Parkinson idiopático, 3 años de evolución
  \item \textbf{MMSE:} 27 (normal)
\end{itemize}

\subsection{Resultados de Pruebas}

\begin{table}[h]
  \centering
  \begin{tabular}{|c|c|c|}
    \hline
    \textbf{Prueba} & \textbf{Resultado} & \textbf{Interpretación} \\
    \hline
    TUG & 9.2 seg & Riesgo Bajo \\
    \hline
    BBS & 48/56 & Riesgo Bajo \\
    \hline
  \end{tabular}
\end{table}

\subsection{Interpretación Clínica}

\begin{itemize}
  \item Riesgo de caída bajo en ambas pruebas
  \item Equilibrio relativamente preservado
  \item Movilidad funcional adecuada
  \item Candidato óptimo para intervención preventiva
\end{itemize}

\subsection{Plan de Intervención}

\begin{enumerate}
  \item Sesiones 2x/semana de 30 minutos
  \item Cueing auditivo a 70 BPM (marcha normal adaptada)
  \item Enfoque en mantención de movilidad y prevención de deterioro
  \item Ejercicios en casa: marcha con metrónomo 10 min/día
\end{enumerate}

\section{Caso 2: Paciente con Riesgo Alto}

\subsection{Presentación}

\begin{itemize}
  \item \textbf{Nombre:} María Rodríguez González
  \item \textbf{Edad:} 75 años
  \item \textbf{Estadio Hoehn-Yahr:} III
  \item \textbf{MMSE:} 25 (límite normal)
  \item \textbf{Antecedentes:} 2 caídas en últimos 3 meses
\end{itemize}

\subsection{Resultados de Pruebas}

\begin{table}[h]
  \centering
  \begin{tabular}{|c|c|c|}
    \hline
    \textbf{Prueba} & \textbf{Resultado} & \textbf{Interpretación} \\
    \hline
    TUG & 22.5 seg & \textbf{RIESGO ALTO} \\
    \hline
    BBS & 38/56 & \textbf{RIESGO ALTO} \\
    \hline
  \end{tabular}
\end{table}

\subsection{Interpretación Clínica}

\begin{advertenciabox}
  \textbf{ALARMA CLÍNICA}
  
  Paciente con riesgo inminente de caídas múltiples. Requiere intervención urgente:
  \begin{itemize}
    \item Evaluación neurológica adicional (posible ajuste medicamentoso)
    \item Valoración oftalmológica (visión)
    \item Valoración de ambiente en el hogar
    \item Posible uso de andador o bastón
  \end{itemize}
\end{advertenciabox}

\subsection{Plan de Intervención Intensivo}

\begin{enumerate}
  \item Sesiones 3x/semana inicialmente
  \item Cueing auditivo a 50 BPM (velocidad reducida)
  \item Supervisión constante (riesgo de caída muy alto)
  \item Considerar derivación a neurólogo para ajuste medicamentoso
  \item Educación familiar sobre factores de riesgo
  \item Seguimiento estricto cada 2 semanas
\end{enumerate}

% CAPÍTULO 7: CONSIDERACIONES FINALES
\chapter{Consideraciones Finales y Mejores Prácticas}

\section{Checklist de Seguridad Antes de Cada Sesión}

\begin{table}[h]
  \centering
  \begin{tabular}{|l|c|}
    \hline
    \textbf{Elemento a Verificar} & \textbf{Confirmado} \\
    \hline
    ¿Paciente tiene tarjeta de citas? & \checkmark \\
    \hline
    ¿MMSE $\geq 24$? & \checkmark \\
    \hline
    ¿Consentimiento informado firmado? & \checkmark \\
    \hline
    ¿Espacio libre de obstáculos? & \checkmark \\
    \hline
    ¿Paciente con calzado seguro? & \checkmark \\
    \hline
    ¿Supervisor/asistente presente? & \checkmark \\
    \hline
    ¿Medicación tomada según protocolo? & \checkmark \\
    \hline
    ¿No hay eventos adversos recientes (caídas, hospitalización)? & \checkmark \\
    \hline
  \end{tabular}
\end{table}

\section{Frecuencia Recomendada de Intervención}

\begin{conceptobox}
  \textbf{Riesgo Bajo (BBS $> 45$):}
  \begin{itemize}
    \item 1-2 sesiones/semana
    \item Duración: 30 minutos
    \item Objetivo: Mantenimiento y prevención
  \end{itemize}
  
  \textbf{Riesgo Moderado (BBS 40-45):}
  \begin{itemize}
    \item 2-3 sesiones/semana
    \item Duración: 30 minutos
    \item Objetivo: Detención de deterioro
  \end{itemize}
  
  \textbf{Riesgo Alto (BBS $< 40$):}
  \begin{itemize}
    \item 3-4 sesiones/semana inicialmente
    \item Duración: 30 minutos con supervisión intensiva
    \item Objetivo: Intervención urgente
    \item Requiere coordinación multidisciplinaria
  \end{itemize}
\end{conceptobox}

\section{Contraindicaciones y Eventos Adversos}

\subsection{Cuándo NO Usar la PWA}

\begin{advertenciabox}
  \begin{itemize}
    \item MMSE $< 24$ (incapacidad cognitiva)
    \item Estadio IV-V Hoehn-Yahr (severidad excesiva)
    \item Insuficiencia cardíaca descompensada
    \item Hipotensión ortostática severa no controlada
    \item Fractura ósea aguda o cirugía reciente
    \item Convulsiones no controladas
    \item Depresión severa con ideación suicida
  \end{itemize}
\end{advertenciabox}

\subsection{Reacciones Adversas Esperadas}

\begin{table}[h]
  \centering
  \begin{tabular}{|p{3cm}|p{4cm}|p{4cm}|}
    \hline
    \textbf{Evento} & \textbf{¿Normal?} & \textbf{Acción} \\
    \hline
    Fatiga leve post-sesión & Sí & Descanso \\
    \hline
    Sudoración & Sí & Monitoreo \\
    \hline
    Mareos leves & Posible & Sentarse, rehidratación \\
    \hline
    Caída durante intervención & No & Detener, evaluar, reportar \\
    \hline
    Dolor torácico & \textbf{EMERGENCIA} & Llamar ambulancia \\
    \hline
  \end{tabular}
\end{table}

\section{Integración en Protocolos Multidisciplinarios}

La PWA es una \textbf{herramienta complementaria}, no un sustituto de:

\begin{itemize}
  \item Seguimiento neurológico regular
  \item Evaluación y ajuste farmacológico
  \item Psicología/psiquiatría (si depresión/ansiedad)
  \item Nutrición especializada
  \item Logopedia (si disfagia)
  \item Terapia ocupacional (si AVD afectadas)
\end{itemize}

El fisioterapeuta debe coordinar con el equipo multidisciplinario para optimizar resultados.

\section{Recomendaciones para Próxima Sesión}

Basándose en resultados actuales:

\begin{enumerate}
  \item \textbf{Si Riesgo Bajo:}
  \begin{itemize}
    \item Incrementar BPM del cueing gradualmente (5 BPM cada 2 sesiones)
    \item Agregar tareas cognitivas durante marcha
    \item Educar sobre actividad física en casa
  \end{itemize}
  
  \item \textbf{Si Riesgo Moderado:}
  \begin{itemize}
    \item Mantener BPM actual
    \item Intensificar trabajo en equilibrio dinámico
    \item Derivar a valoración de ambiente en hogar
  \end{itemize}
  
  \item \textbf{Si Riesgo Alto:}
  \begin{itemize}
    \item Mantener BPM bajo (40-50)
    \item Máxima supervisión
    \item Contacto inmediato con médico tratante para revisión de medicación
  \end{itemize}
\end{enumerate}

\appendix

\chapter{Glosario de Términos Clínicos}

\begin{description}
  \item[\textbf{Bradicinesia}] Lentitud patológica de movimientos voluntarios
  
  \item[\textbf{Congelamiento (Freezing)}] Episodios súbitos e involuntarios de inmovilidad, particularmente durante la marcha
  
  \item[\textbf{Cueing}] Técnica de proporcionar pistas externas (auditivas, visuales, táctiles) para facilitar movimiento automático
  
  \item[\textbf{Escala Hoehn-Yahr}] Sistema de clasificación de severidad del Parkinson (Estadios 0-V)
  
  \item[\textbf{Fenómeno ``Off''}] Período de pérdida de eficacia de medicación antiparkinsonian
  
  \item[\textbf{Freezing of Gait}] Pérdida temporal e involuntaria de movimiento durante la marcha
  
  \item[\textbf{MMSE}] Mini Mental State Examination, prueba de screening cognitivo
  
  \item[\textbf{Neuroplasticidad}] Capacidad del sistema nervioso de reorganizarse y adaptarse
  
  \item[\textbf{Propiocepción}] Sentido de posición del cuerpo en el espacio
  
  \item[\textbf{Rigidez}] Resistencia al movimiento pasivo de una articulación
  
  \item[\textbf{Temblor en Reposo}] Oscilación involuntaria característica del Parkinson en posición estática
  
  \item[\textbf{UPDRS}] Unified Parkinson's Disease Rating Scale, escala de evaluación estándar internacional
\end{description}

\chapter{Referencias y Bibliografía Recomendada}

\begin{enumerate}
  \item Hoehn, M. M., \& Yahr, M. D. (1967). Parkinsonism: onset, progression and mortality. \textit{Neurology}, 17(5), 427-442.
  
  \item Berg, K. O., Wood-Dauphinee, S. L., Williams, J. I., \& Maki, B. (1992). Measuring balance in the elderly: validation of an instrument. \textit{Canadian Journal of Public Health}, 83, 7-11.
  
  \item Podsiadlo, D., \& Richardson, S. (1991). The timed ``Up and Go'': a test of basic functional mobility for frail elderly persons. \textit{Journal of the American Geriatrics Society}, 39(2), 142-148.
  
  \item Walshe, E. A., Pereira, B. P., \& Marras, C. (2015). Premotor Parkinson's disease: identifying early Parkinson's disease. \textit{Parkinson's Disease}, 2015, 364058.
  
  \item Almeida, L. R. S., Sherrington, C., Allen, N. E., \& Canning, C. G. (2015). Rehabilitation aimed at improving balance, functional mobility, and psychosocial functioning reduces falls in older adults with Parkinson's disease. \textit{Parkinson's Disease}, 2015, 709168.
\end{enumerate}

\end{document}
